% Transcription — fonds Alexandre Grothendieck, shelfmark 58, batch 3 (pages 41-60).
% Established from the facsimile archives/batches/58/batch-03.pdf (PDF page k =
% archivists' page 40+k, no cover sheet), cut from a copy of 58.pdf fetched
% through the project's own relay
% (https://grothendieck-relay.onrender.com/source/58.pdf); tiles in
% archives/tiles/58/p41-p60, re-cropped at 330 dpi with pdftoppm. Per the
% skill transcribe-grothendieck. The folder is sixty-six pages in four batches
% (1-20, 21-40, 41-60, 61-66), transcribed in parallel; this is the third,
% written after batch 2 (whose notation it keeps: (LG)_n, LG faible/stricte,
% X_{/Y}) and with only a glance at the header of batch 4.
%
% Pass: Opus 5.5 (claude-opus-5-5), 2026-09-26 — first pass, unchecked against the pages by a human.
%
% Pages skipped: 44 and 54 — blank cream sheets. 46, 48 — typed leaves of a
% Bourbaki draft (« II - 11 -, II - 10 -, n° 338 »: projective modules of rank
% 1, fractional ideals, class group), nothing in his hand. 50, 52 — typed
% leaves scanned upside down (another typescript, n° 3??), nothing in his
% hand; skipped by the settled rule for typed versos. Page 53 is a tan cover
% inscribed top right « Applications aux morphismes [struck word] projectifs
% et aux morphismes propres »; it takes no \page marker and its words become
% the section heading, with a \note.
%
% Pages summarised: none. Pages transcribed: 41, 42, 43 (pencil on grey
% paper), 45, 47, 49, 51 (ink on the blank sides of typed leaves, whose
% typescript shows through), 55, 57, 58, 59 (blue ink), 56 and 60 (two copies
% of one typed page — the end of a statement on the existence of a coherent
% sheaf with given formal completion, « Corollaire 3.2. Sous les conditions
% de 3.1. (» — carrying his ink (56) and pencil (60) calculations and
% sketches; the typescript is described in one \note each and his additions
% are given in full). Page 45 and 49 are palimpsests with long diagonal
% margin notes and zigzag cancellations; the formulas carry, much prose does
% not. Pages 55-59 are written fairly and read almost end to end.
%
% His pagination: none visible. His numbered runs: p. 41 continues « 3)
% Théorème de permanence » of p. 40 (the cancelled Th. 2 ends at the top of
% p. 41, then a new unnumbered « Th. » with (i), (ii), (iii)); Corollaire 2
% (p. 42) precedes Corollaire 1 (p. 42, (X quasi-projectif)); « Remarques »
% (p. 43) refer to « le Th. et les deux corollaires » and « Cor. 2 (ii) ».
% P. 45: an unnumbered « Th. » on a local ring A, t in rA; p. 49 cites « cor.
% 3 au th. de comparaison affine » and conditions a) b) c'), T_1 (not on these
% leaves). Pp. 55-59: questions a), a'), b), c) (p. 55), then circled
% paragraphs (1°) (p. 55), (2°) (p. 58), (3°) (p. 59), with circled labels
% for groups of equivalent conditions: « * » (p. 55), « B_n » (p. 57), « A'_i »
% (p. 58), « A^f_n » (p. 59). Batch 4's p. 61 continues with groups C'_i, C_n.
%
% What the batch is about. Pp. 41-43: permanence of the Lefschetz-Grauert
% conditions under f : X' -> X projective (and flat): (i) (LG)_n gives
% H^i(X',F') -> H^i(X'_{/Y'},F'_{/Y'}) iso for i <= n, mono for i = n+1,
% for F' coherent, flat, with support proper over X; (ii) LG gives full
% faithfulness of F' -> F'_{/Y'} on flat coherent Modules (margin: Hom(F',G')
% iso when G' is flat); (iii) LG strict gives algebraisation of flat formal
% Modules on a neighbourhood; Cor. 2: the same for projective flat schemes
% over X versus formal schemes over X_{/Y}; Cor. 1: (LG)_n, LG, LG strict
% pass to (X', f^{-1}(Y)) for X' projective flat over X quasi-projective;
% Remarques on the necessity of flatness and properness. Pp. 45-51: a local
% theorem — A local noetherian, quotient of a regular ring, t in rA,
% X = Spec A, Y = V(t), X' = X - {x}, Y' = Y cap X', Phi the finite closed
% subsets of X' outside Y'; the functor F -> F^ from coherent sheaves on X'
% to coherent sheaves on X^' has kernel C_Phi(X') and induces a fully
% faithful rho on the quotient category, via Hom(F,G) = Hom(F^,G^) when G has
% depth >= 2 off Y', a dévissage with a four-term exact sequence, the
% t^n-torsion kernel and a diagram of Gamma's and H^1's; pp. 49-51 prove that
% a formal sheaf satisfying a), b), c') comes from an F' of depth >= 2 at the
% points of T_1 (« on a gagné »). Pp. 55-60 (after the cover « Applications
% aux morphismes projectifs et aux morphismes propres »): Y a hyperplane
% section of X = P^r_k, U a neighbourhood of Y, Z = X - U; the questions of
% vanishing and finiteness of H^i(U, F(n)) and of H^i(U,F) -> H^i(X^,F^)
% being iso; the local cohomology sequence H^i_Z(X,F(-n)) -> H^i(X,F(-n)) ->
% H^i(U,F(-n)) -> H^{i+1}_Z, dualised by E^j = Ext^j(F, omega_X); chains of
% equivalences (Supp E^{r-i}, points isolated, prof F_x > n, R^i g_*(F)
% coherent, H^i_z(F_z) of finite length) for alpha_i injective/bijective.
%
% Notation, kept from batch 2: X_{/Y}, F_{/Y}, \hat{X}, \hat{F}; (\mathrm{LG})_n
% and (LG), LG faible/stricte as he writes them; \mathfrak{X} for his
% gothic X (= X_{/Y}), \mathfrak{X}' for a formal scheme over it, \mathcal{F}'
% for his script F (a formal Module); \Phi for his looped capital phi (as in
% batch 1); \mathfrak{r}A for the radical; C(X'), C_\Phi(X') for his script C;
% \underline{\mathrm{Ext}}, \underline{H} for the underlined sheaf functors.
% On pp. 55-60 the twist is m (his long wavy letter, as batch 4 reads it on
% p. 61: F(m), F(-m), m grand) and n the bound on i (i <= n, B_n, A^f_n).
%
% Readings to check first: p. 41 « plat sur Y » (X expected) and the margin
% on Hom(F',G'); the margin of p. 42; pp. 43 bottom (the boxed insertion);
% p. 45 second line and the long diagonal margin notes; p. 47 « t^n » and the
% letter G' (perhaps C'), the arrows of the diagram and its left label; the
% zigzag-cancelled middle of p. 49; p. 55 « injectif » written over another
% word in the first group; p. 57 « sur S » and the conclusion lines; the
% subscripts of the circled labels (p. 57 B_n, p. 59 A^f_n).

\documentclass[11pt,a4paper]{article}
% Le préambule commun à toutes les transcriptions du fonds Grothendieck.
%
% Il tient en une idée : **l'incertitude se compose, elle ne se résout pas.**
% Une transcription qui lisse un mot douteux a détruit la seule chose pour
% laquelle on la faisait — un lecteur doit pouvoir savoir, à chaque endroit,
% ce qui était sur le feuillet et ce qui a été ajouté. Les six macros qui
% suivent sont donc l'appareil critique, et non de la décoration.
%
% Les mêmes commandes sont reconnues par `scripts/render.mjs`, qui produit la
% vue de lecture du site. Une macro ajoutée ici doit l'être aussi là-bas,
% faute de quoi le rendu échouera bruyamment — ce qui est voulu : un rendu
% silencieusement faux est pire qu'un rendu absent.

\ProvidesPackage{grothendieck}[2026/08/08 Transcriptions du fonds Grothendieck]

\usepackage{amsmath,amssymb,amsthm}
\usepackage{mathrsfs}
\usepackage{stmaryrd}
% Les diagrammes commutatifs sont partout dans ce fonds : c'est la langue
% même de la géométrie algébrique de Grothendieck, et une transcription qui
% les rendrait en prose ne transcrirait rien.
\usepackage{tikz-cd}
% Français par défaut — c'est la langue des feuillets — mais l'anglais est
% chargé aussi : la traduction et le résumé sont des documents anglais, et la
% typographie française (espaces avant les deux-points) y serait une faute.
% Ces documents basculent par \selectlanguage{english} en tête de corps.
\usepackage[english,french]{babel}
\usepackage{fontspec}
\usepackage{geometry}
\geometry{a4paper,margin=25mm,marginparwidth=28mm}
\usepackage{marginnote}
\usepackage{xcolor}
% ulem, not soul. `soul`'s \st and \ul are built on a character-by-character
% scan that XeLaTeX's font machinery breaks: the document fails with an
% undefined control sequence rather than degrading. ulem does the same two jobs
% and is Unicode-engine safe. `normalem` stops it hijacking \emph, which the
% transcriptions use for Grothendieck's own emphasis.
\usepackage[normalem]{ulem}
\usepackage{hyperref}
\hypersetup{colorlinks=true,linkcolor=black,urlcolor=black,pdfborder={0 0 0}}

% --- Métadonnées du lot -----------------------------------------------------
% Reprises telles quelles par le script de rendu, qui les affiche en tête de
% la vue de lecture. Elles fixent ce que le fichier prétend couvrir : sans
% elles, une transcription flotte sans point d'attache dans un fonds de seize
% mille feuillets.

\newcommand{\folder}[1]{\gdef\grothFolder{#1}}
\newcommand{\batch}[1]{\gdef\grothBatch{#1}}
\newcommand{\leaves}[2]{\gdef\grothFirst{#1}\gdef\grothLast{#2}}
\newcommand{\foldertitle}[1]{\gdef\grothFoldertitle{#1}}
\gdef\grothFolder{?}\gdef\grothBatch{?}\gdef\grothFirst{?}\gdef\grothLast{?}\gdef\grothFoldertitle{}

% --- Le résumé --------------------------------------------------------------
%
% Il ouvre la lecture modernisée, sous le titre « Résumé », et il est composé
% en petit. Ce n'est pas un ornement : c'est le seul passage du document qui
% ne soit pas des mathématiques, et un lecteur doit pouvoir le reconnaître
% avant de l'avoir lu — puis le sauter s'il connaît déjà le sujet, ou s'y
% arrêter s'il ne le connaît pas. Le filet qui le suit marque la reprise du
% corps.

\newenvironment{resume}
  {\small\setlength{\parskip}{0.45em}}
  {\par\medskip\hrule\medskip}

% --- Mots-clés ---------------------------------------------------------------
%
% En anglais, en clôture du résumé de la lecture modernisée : le vocabulaire
% moderne sous lequel chercher ce que les feuillets construisent. Le manifeste
% du site (`npm run manifest`) les extrait pour étiqueter la cote entière —
% cette ligne est la source unique des tags, il n'y a pas de fichier à côté.
\newcommand{\keywords}[1]{\par\smallskip\noindent
  {\footnotesize\textsc{Keywords} --- \textit{#1}}\par}

% --- L'appareil critique ----------------------------------------------------

% Le numéro de feuillet, en marge. C'est l'ancre qui permet de régler un
% désaccord sur une formule en regardant *un* feuillet plutôt qu'une chemise
% de six cents. Le site s'en sert aussi pour tourner les pages du fac-similé
% quand on fait défiler la transcription.
\newcommand{\leaf}[1]{%
  \marginnote{\footnotesize\color{gray!70}\textsf{#1}}%
  \label{leaf:#1}\ignorespaces}

% Illisible : on ne devine pas. Un mot inventé qui se lit comme les autres est
% le pire résultat possible.
\newcommand{\ill}{\textcolor{red!60!black}{[\,\ldots\,]}}

% Lecture douteuse : le mot est donné, et signalé comme incertain.
\newcommand{\uncertain}[1]{\uline{#1}}

% Ajout de l'éditeur — une lettre manquante, un mot sous-entendu.
\newcommand{\add}[1]{\textcolor{blue!50!black}{[#1]}}

% Biffé par Grothendieck lui-même. Ce qu'il a rayé fait partie du document :
% une rature dit souvent où la pensée a changé de direction.
\newcommand{\struck}[1]{\sout{#1}}

% Note du transcripteur, sur l'état matériel du feuillet.
\newcommand{\note}[1]{\par\smallskip{\small\color{gray!60!black}\textit{[#1]}}\par\smallskip}

% Note marginale de Grothendieck, distincte du corps du texte.
\newcommand{\marginal}[1]{\par\smallskip\noindent\hspace{1em}%
  {\small\color{gray!60!black}#1}\par\smallskip}

% --- Titre ------------------------------------------------------------------

\AtBeginDocument{%
  \begin{center}
    {\large\bfseries Fonds Alexandre Grothendieck --- cote n\textsuperscript{o}~\grothFolder}\\[2pt]
    {\itshape\grothFoldertitle}\\[4pt]
    {\small feuillets \grothFirst--\grothLast\ (lot \grothBatch)}\\[2pt]
    {\footnotesize\color{gray!60!black}Transcription établie d'après le fac-similé numérisé
     par l'Université de Montpellier.\\
     Les lectures incertaines sont soulignées, les illisibles notées [\,\ldots\,],
     les ajouts entre crochets.}
  \end{center}
  \vspace{1em}}

\endinput


\folder{58}
\batch{3}
\pages{41}{60}
\dating{[à partir de 1964]}
\watermark{Édition de démonstration}
\foldertitle{Théorie de " Lefschetz-Grauert ". Applications en π₁, à Pic etc… : notes manuscrites (s.d.)}

\begin{document}

\page{41}
\struck{pour tout $f\colon X' \to X$ \ill{} projectif et plat, posant $Y' = f^{-1}(Y)$, $(X', Y')$ satisfait (LG) [resp.\ LG strict].}
\note{Fin de l'énoncé du Th.~2 de la page 40, lui aussi annulé : les trois premières lignes de la page sont barrées de deux longs traits obliques et fermées par un trait horizontal.}

\emph{Th.} Soient $X$ \struck{préschéma} \struck{loc.} noeth.\ [\emph{quasi-projectif}], $Y$ partie fermée, $f\colon X' \to X$ un \uncertain{morphisme} \uncertain{projectif}, $Y' = f^{-1}(Y)$. Soit $F'$ un faisceau coh.\ sur $X'$, à support propre sur $X$, \emph{plat} sur $Y$, \uncertain{considérons} les homom.
\[
  (*) \qquad H^i(X', F') \longrightarrow H^i(X'_{/Y'}, F'_{/Y'})
\]
\struck{Sous ces conditions,}
\note{« quasi-projectif », doublement souligné, est ajouté en haut à droite et relié par un crochet ; l'hypothèse « à support propre sur $X$ » est une insertion encadrée. On lit $f^{-1}(Y)$, peut-être suivi d'un prime. On attendrait « plat sur $X$ » : la lettre se lit $Y$.}

(i) Si $(X, Y)$ satisfait $(\mathrm{LG})_n$, alors $(*)$ est un isom.\ pour $i \leq n$, un \uncertain{monom.}\ pour $i = n+1$.

(ii) Si $(X, Y)$ satisfait (LG), alors pour tout voisinage ouvert $U$ de $Y$, posant $U' = f^{-1}(U)$, \struck{l'homom.}\ le foncteur \struck{\ill{}} $F' \mapsto F'_{/Y'}$ des Modules \ill{} coh.\ sur $U'$ plats $/Y$ [et : à support propre $/X$], dans les Modules coh.\ sur $X'_{/Y'}$ [à support propre sur $\mathfrak{X}$], plats sur $X_{/Y} = \mathfrak{X}$, est pleinement fidèle.
\marginal{$\operatorname{Hom}(F', G') \xrightarrow{\ \sim\ } \operatorname{Hom}(F'_{/Y'}, G'_{/Y'})$ pourvu que $G'$ soit plat sur $X$ [$F'$ \ill{}] [\ill{} \ill{} \ill{} \ill{}\ldots]}
\note{Note marginale gauche, écrite en oblique et reliée à (ii) ; son $G'$ est doublement souligné.}

(iii) Si $(X, Y)$ satisfait (LG) strict, alors pour tout Module coh.\ $\mathcal{F}'$ sur $X'_{/Y'}$, plat rel.\ $X_{/Y}$, il

\page{42}
existe un voisinage ouvert $U$ de $Y$, \struck{\ill{}} un module $F'$ sur $U' = f^{-1}(U)$, plat $/U$, et un isom.
\[
  F'_{/Y'} \simeq \mathcal{F}' .
\]

\emph{Corollaire 2} Sous les conditions \ill{} \ill{} \ill{} $(X, Y)$, [(i) Si $(X, Y)$ satisfait (LG),] le foncteur $X' \mapsto X' \times_X X_{/Y}$ de la catégorie des \struck{pré}schémas \struck{\ill{}} projectifs et plats sur $X$ dans la catégorie des schémas formels projectifs et plats sur $X_{/Y}$ est pleinement fidèle. Kif Kif en remplaçant $X$ par un voisinage ouvert $U$ de $Y$.
\marginal{\ill{} pour les \uncertain{préschémas} \emph{\ill{}} ; \ill{} \uncertain{faisceaux} coh.\ \emph{plats} et \emph{fidèles} \ill{}}
\note{« (i) Si $(X, Y)$ satisfait (LG) » est écrit en interligne et relié par un crochet. Note marginale gauche en oblique, en partie soulignée, presque illisible.}

(ii) Si $(X, Y)$ satisfait (LG strict), pour \struck{tout} $\mathfrak{X}'$ [projectif et de type fini] (\struck{\ill{}} $\mathfrak{X}'/X_{/Y}$) plat sur $X_{/Y}$, il existe un voisinage ouvert $U$ de $Y$, \struck{\ill{}} un $U'$ projectif \struck{de type fini} et plat sur $U$, et un $\mathfrak{X}$-isom.
\[
  U'_{/Y'} \simeq \mathfrak{X}' \quad (\text{où } Y' = f^{-1}(Y)).
\]

\emph{Corollaire 1} [($X$ quasi-projectif)] (i) Si $(X, Y)$ satisfait $(\mathrm{LG})_n$ \struck{alors} [resp.\ LG, resp.\ LG strict], alors pour tout $X'$ projectif et plat sur $X$, posant \struck{$X$} $Y' = f^{-1}(Y)$, $(X', Y')$ satisfait $(\mathrm{LG})_n$ [resp.\ (LG), resp.\ (LG) strict].

\struck{(ii) Si $(X, Y)$}
\note{« ($X$ quasi-projectif) » est encadré au-dessus de l'énoncé ; les crochets « [resp.\ LG, resp.\ LG strict] » et « [resp.\ (LG), resp.\ (LG) strict] » sont ajoutés à l'encre brune. Le Corollaire~1 est écrit après le Corollaire~2 sur la page.}

\page{43}
\emph{Remarques} On voit facilement par des exemples dans le Th.\ et les deux corollaires, que les conditions de platitude et de \emph{propreté} sont essentielles. Cependant, il \struck{semble} est possible que l'hypothèse de quasi-projectivité sur $X$ [et $X'$] \struck{et celle de} peut être omise, \struck{et celle de propreté dans les corollaires, dans le th., et les} \ill{} \ill{} \ill{} \ill{} en plus dans [\uncertain{partie des}] Cor.~2 (ii) [Existence d'un schéma algébrique donnant un schéma formel].
\note{« et $X'$ » est ajouté en interligne. Une ligne et demie est biffée ; sous elle un petit cadre relié par une boucle porte « partie des » (lecture incertaine), qui s'insère avant « Cor.~2 (ii) ». Le reste de la page est blanc.}

\page{45}
\emph{Th.} $A$ anneau local noeth.\ [quotient d'un régulier]. $t \in \mathfrak{r}A$, $A$ \struck{\ill{}} [\ill{}] \ill{} la top.\ $t$-adique \ill{} \ill{}
$X = \operatorname{Spec}(A)$, $Y = V(t)$, $x$ \ill{} $\mathfrak{m} = \mathfrak{r}A$, $X' = X - \lbrace x \rbrace$, $Y' = Y - \lbrace x \rbrace = Y \cap X'$,
\[
  \hat{X} = X_{/Y}, \qquad \hat{X}' = X'_{/Y'},
\]
$\Phi$ l'ens.\ des parties fermées de $X'$ \ill{} \ill{} \ill{} $Y'$, i.e.\ \emph{finies}, $\mathcal{C}(X')$ la catégorie des faisceaux coh.\ sur $X'$, $\mathcal{C}_\Phi(X')$ la sous-catégorie des faisceaux \struck{\ill{}} à supports $\in \Phi$. Considérons le foncteur [\emph{exact}]
\[
  F \mapsto \hat{F} : \mathcal{C}(X') \longrightarrow \mathcal{C}(\hat{X}') .
\]
\note{Encre noire sur le côté blanc d'un feuillet dactylographié dont le texte transparaît. La deuxième ligne est surchargée ; « $t$-adique » est d'une lecture incertaine. Une longue insertion écrite en diagonale à gauche, reliée par des flèches à « $\Phi$ » et à « $\mathcal{C}_\Phi(X')$ », est illisible.}

Alors le foncteur est exact, a noyau $\mathcal{C}_\Phi(X')$, [donc définit un foncteur [conservatif exact, fidèle]
\[
  \rho : \mathcal{C}(X') / \mathcal{C}_\Phi(X') \longrightarrow \mathcal{C}(\hat{X}') .]
\]
De plus, $\rho$ est pleinement fidèle / \struck{\ill{}}, et définit une \emph{\uncertain{équivalence}} de $[\mathcal{C}(X')/\mathcal{C}_\Phi(X')]'$ \uncertain{avec} la catégorie des Modules coh.\ sur $\hat{X}'$ [\emph{\uncertain{Dominés}}] qui sont restriction de Modules coh.\ sur $\hat{X}$, et \ill{} \ill{} \struck{\ill{} \ill{} \ill{} \ill{}}]
\marginal{\struck{[\ill{} \ill{} \ill{} \ill{} \ill{} \ill{} \ill{} \ill{} \ill{}] [\ill{} \ill{} \ill{} \ill{} \ill{} 2 \ill{} 3]} \ill{} \ill{} \ill{} \ill{}}
\note{La marge gauche porte, sur toute la hauteur du milieu de la page, cinq ou six lignes obliques serrées, la plupart barrées d'un trait, reliées au texte par des flèches ; on n'y distingue que quelques chiffres. La fin de l'alinéa est surchargée, avec un ajout encerclé sous la ligne, illisible.}

Il faut surtout prouver que le foncteur $\rho$ est \emph{pleinement fidèle}, on montre ceci : si $F$, $G$ \struck{sont} coh.\ sur $X'$, \struck{\ill{}} $G$ de prof $\geq 2$ en les \struck{$x \in \Phi$ \ill{}} points \ill{} de $X' - Y'$, \struck{\ill{}} alors
\[
  \operatorname{Hom}(F, G) \simeq \operatorname{Hom}(\hat{F}, \hat{G}) .
\]
Pour ce dernier point, en considérant $H = \underline{\mathrm{Hom}}(F, G)$, on est ramené à prouver que
\[
  \Gamma F \simeq \Gamma \hat{F}
\]
si $F$ est de prof $\geq 2$ en les \uncertain{pts} [fermés] de $X' - Y$.
\note{Dans la formule $\Gamma F \simeq \Gamma \hat{F}$, écrite en surcharge, les lettres sont récrites sur d'autres ; peut-être $\Gamma(X', F) \simeq \Gamma(\hat{X}', \hat{F})$.}

\page{47}
Pour le voir, on est ramené \ill{} un \ill{} dévissage facile, [utilisant l'existence d'une suite exacte
\[
  0 \to P' \to F' \to \bar{F}' \to Q' \to 0
\]
où $\operatorname{Supp} P' \subset Y'$, $\operatorname{Supp} Q' \subset Y'$, $\dim \operatorname{supp} P' \leq 0$, $\dim \operatorname{supp} Q' \leq 1$, et $\operatorname{prof} \bar{F}'_x \geq$ \ill{} en \uncertain{tout} pt fermé de $X'$ \ill{} \ill{} $X' - Y'$ ; \ill{} dévissage qui est une question purement locale sur \ill{} \ill{} \ill{} \emph{\ill{}} $X'$] \ill{} \ill{} \ill{} \ill{}.
\note{Encre noire sur le côté blanc d'un feuillet dactylographié. La barre de $\bar{F}'$ et les bornes $\leq 0$, $\leq 1$ sont lues d'après le tracé.}

$\operatorname{prof} F_x \geq 2$. \emph{D'ailleurs}, \ill{} \ill{} pour \emph{tout} $x \in X'$. Soit d'ailleurs $K'$ le noyau de la multiplication par \uncertain{$t^n$} \struck{\ill{}} (\ill{} $n$ \ill{}) \ill{}
\[
  0 \to K' \to F' \to G' \to 0
\]
avec $G'$ \struck{\ill{}} $t$-régulier. \ill{}.
\note{La lettre lue $G'$ pourrait être $C'$. En marge gauche, deux signes encerclés (un cercle traversé de traits, avec un point) marquent les deux alinéas ; leur sens n'est pas clair.}

\[
\begin{array}{ccccccc}
  0 \to \Gamma K' & \longrightarrow & \Gamma F' & \longrightarrow & \Gamma G' & \longrightarrow & H^1 K' \\
  \downarrow \wr & & \downarrow & & \downarrow \wr\,? & & \wr \\
  0 \to \Gamma \hat{K}' & \longrightarrow & \Gamma \hat{F}' & \longrightarrow & \Gamma \hat{G}' & \longrightarrow & H^1 \hat{K}'
\end{array}
\]
\note{À gauche de la seconde ligne du diagramme, les signes « $A$ $\hat{A}_{p}$ » (lecture incertaine) ; le « ? » après la troisième flèche verticale est de lui.}

Le \ill{} \ill{} \ill{} \ill{} \ill{} \ill{} \ill{} en \ill{} \ill{}.
\note{Un trait horizontal clôt le texte ; le bas de la page est vide.}

\page{49}
\struck{\ill{} \ill{} \uncertain{établissons}}

\emph{Avec} cor.~3 au th.\ de comparaison affine supposons que $\mathcal{F}'$ satisfasse a) b) c') [c') \ill{} \ill{} $\mathcal{F}'$ s'étende en un $\mathcal{F}$ coh.\ sur $\hat{X}$], et disons que $\mathcal{F}'$ provient d'un $F'$ satisfaisant à la condition $\operatorname{prof} F'_x \geq 2$ si $x \in T_1$. En effet, il provient d'ailleurs d'un $F'$ coh.\ \struck{\ill{}} [\ill{}] $t$-régulier, et un $F'$ $t$-régulier, il faut \ill{} \ill{} \ill{} \ill{} condition $F'$. \struck{Soit}
\note{Les conditions a), b), c') et l'ensemble $T_1$ ne sont pas définis sur ces feuillets.}

\struck{$Z'$ l'ens.\ des pts de $X'$ tels que $\operatorname{prof}_{x'}(F') \leq 1$, [c'est une partie \uncertain{constructible} \ill{} \ill{} \ill{}] \ill{} $X' - Y'$ \ill{} \ill{} \ill{} \ill{}, \ill{} \ill{} \ill{} $x' \in \operatorname{Ass} F'$, \ill{} \ill{} ; $x' \notin T_1$ \ill{} $\overline{x}' \cap Y' \neq \emptyset$, \ill{} \ill{} \ill{} \ill{} définition $F'$ \ill{} \ill{} \ill{} \ill{} \ill{} \ill{} des \ill{} \ill{} $U$ \ill{} \ill{} \ill{} $x'$ \ill{} $\operatorname{prof}_{x'}(F') \geq 2$ \ill{} \ill{} \emph{fermés} de $Y'$,}
\marginal{$\overline{x}'$ \ill{} fermé dans $X'$}
\marginal{Je suppose $X$ \uncertain{localisé} autour de $T$, et $X - T = X'$ \struck{\ill{}} \ill{} « \ill{} » \ill{}}
\note{Tout ce passage, dans un cadre, est annulé par un long trait en zigzag ; on n'en donne que les formules et quelques mots. Les deux notes marginales gauches y sont reliées, la première dans une bulle.}

\ill{} \ill{} \ill{}. $x' \in U$, c'est donc \ill{} \ill{} \ill{} $U$ \ill{} \ill{} $Y$, donc $U$ est un voisinage de $Y'$ dans $X'$. \struck{\ill{} \ill{} $x' \in X' - U$ \ill{} \ill{}} Les pts de $X' - U$ sont donc \emph{fermés} dans $X'$ [et non \ill{} \ill{} \ill{} $T_1$] --- donc sont \ill{} partie \ill{} finie. Soit $x' \in X' - U$, et \struck{$x'_1 \in \operatorname{Ass} F'$ tel} $x'_1 \in U$ tel que $x'_1 \in \operatorname{Ass} F'$, $\overline{x'_1} \ni x'$, je dis que \ill{} \ill{} \ill{} \ill{} \ill{} \ill{} $x'_1$ \ill{} $x'$ : \ill{} en effet, $\overline{x'_1}$ serait de dim~1,

\page{51}
donc en les pts fermés $z'$ de $\overline{x'_1} \cap U$, on aurait $\operatorname{prof} F'_{z'} \leq 1$, ce qui est absurde. Donc si $i\colon U \to X'$ est l'inclusion, $F'_1 = i_*(F' \mid U)$ est cohérent, et \uncertain{évidemment} de $\operatorname{prof} \geq 2$ en $X' - U$, d'autre part il est de $\operatorname{prof} \geq 2$ aux pts fermés de \struck{$X'$} $U$, donc $F'_1$ est de $\operatorname{prof} \geq 2$ en \emph{tt} pt fermé de $X'$ ; on a gagné.
\note{Encre noire sur le côté blanc d'un feuillet dactylographié ; le reste de la page est vide.}

\section{Applications aux morphismes projectifs et aux morphismes propres}
\note{Titre écrit par Grothendieck en haut à droite de la couverture brune (page 53), souligné ; un mot biffé et hachuré précède « projectifs ».}

\page{55}
$Y \subset \mathbf{P}^r_k = X$ (section hyperplane), $U$ voisinage de $Y$, $Z = X - U$,

$F$ cohérent sur $X$, \ill{} \ill{} \ill{}

a) Condition pour que $H^i(U, F(-m)) = 0$ pour $m$ grand ?

a') Condition pour que $\coprod_m H^i(U, F(m))$ de type fini sur $S$

b) Condition pour que $H^i(U, F(m)) = 0$ pour $m$ \uncertain{grand} ($i > 0$ donné)

c) Condition pour que
\[
  H^i(U, F) \longrightarrow H^i(\hat{X}, \hat{F})
\]
soit un isomorphisme, injectif ?
\note{« c) » est écrit sur un autre signe ; dans $H^i(\hat{X}, \hat{F})$, le $\hat{X}$ est récrit sur un $X$.}

(1°) Considérons la suite exacte (où $E^j = \underline{\mathrm{Ext}}^j_{\mathcal{O}_X}(F, \omega_X)$)
\note{« 1° » est encerclé ; de même « 2° » et « 3° » aux pages 58 et 59.}
\[
  \begin{array}{l}
  H^i_Z(X, F(-m)) \xrightarrow{\alpha_i} H^i(X, F(-m)) \to H^i(U, F(-m)) \\
  \qquad \to H^{i+1}_Z(X, F(-m)) \xrightarrow{\alpha_{i+1}} H^{i+1}(X, F(-m))
  \end{array}
\]
et, en dessous, les duaux :
\[
  \prod_{z \in Z} E^{r-i}_z(m)^{\wedge} \longleftarrow \Gamma E^{r-i}(m)
  \qquad\qquad
  \prod_{z \in Z} E^{r-i-1}(m)_z^{\wedge} \longleftarrow \Gamma E^{r-i-1}(m)
\]
\note{Les deux flèches de la seconde ligne sont écrites sous $\alpha_i$ et $\alpha_{i+1}$ respectivement (dualité de Serre) ; « et, en dessous, les duaux » est de l'éditeur, non de lui. Les arguments $F(-m)$ sont récrits sur $F$.}

Conditions équivalentes
\[
\left.
\begin{array}{l}
  \alpha_i \text{ injectif pour } m \text{ grand} \\
  \Updownarrow \\
  z \in Z \Rightarrow z \text{ isolé dans } (\operatorname{Supp} E^{r-i} \cup \lbrace z \rbrace) \\
  \Updownarrow \\
  H^{i-1}(X, F(-m)) \to H^{i-1}(U, F(-m)) \text{ surjectif pour } m \text{ grand}
\end{array}
\right.
\]
\note{« injectif » est écrit sur un premier mot (« surjectif » ?) ; à la suite, un crochet biffé « [\ill{} pour tout \ill{}] ». En formule, les mots sont donnés sans l'appareil ; la lecture « injectif » reste incertaine.}

\[
(*) \quad
\left.
\begin{array}{l}
  \alpha_i \text{ injectif pour } m \text{ grand},\ i \leq n \\
  \Updownarrow \\
  \exists \text{ voisinage ouvert } V \text{ de } Z, \text{ tel que } F \text{ soit de profondeur } \ldots \text{ sur } U \cap V \\
  \Updownarrow \\
  H^i(X, F(-m)) \to H^i(U, F(-m)) \text{ surjectif pour } m \text{ grand},\ i < n
\end{array}
\right.
\]
\note{La borne de la profondeur, sur la deuxième ligne, n'est pas écrite ou ne se lit pas (on la rend par des points) ; après « $i \leq n$ », une flèche $\Rightarrow$ est biffée. L'astérisque est en marge gauche de l'accolade.}

\page{56}
\note{La moitié supérieure de la page est un feuillet dactylographié : la fin d'un énoncé où $F_0$ est le Module induit sur $X_0$, sous les hypothèses a) $t_s$ est $\mathcal{O}_{X_s}$-régulière et $F_s$-régulière, b) $F$ et $\mathcal{O}_X$ plats sur $S$, c) $F_{0\,s}$ et $\mathcal{O}_{X_{0\,s}}$ de profondeur $\geq 2$ aux points fermés de $X_s$, avec pour conclusion l'existence d'un Module cohérent sur $X$ dont le complété formel le long de $X_0$ est isomorphe au Module formel donné ; puis « Cet énoncé va résulter du suivant : Corollaire~3.2. Sous les conditions de 3.1. ( », qui s'interrompt. Une retouche à la machine corrige « et $\mathcal{O}_X$ ». Ce qui suit est de sa main, à l'encre.}

\[
  \coprod_{\substack{p \geq 0 \\ q \geq 0}} F_0(q - p)
  \qquad\qquad
  \mathcal{M} + \mathcal{M}(1) + \cdots + \mathcal{M}(-p)
\]
\[
  \coprod_{q} S_q
\]
\note{Sous l'indice du second coproduit, un « $0$, » précède le $q$. Suivent deux croquis : à gauche, dans un quadrant d'axes, un triangle hachuré coupé par une diagonale et une courbe sinueuse, avec les légendes « $\sigma F(n)$ » et « $\mathcal{F}(n)$ » ; à droite, un quadrant hachuré de droites parallèles, légendé en haut « \ill{} $q$ » et à droite « $p$ » et « \uncertain{filtre} ».}
\[
  H^i(X, F) \longrightarrow H^i(\hat{X}, \hat{F})
\]

\page{57}
\[
\left.
\begin{array}{l}
  \alpha_i \text{ bijectif pour } m \text{ grand} \\
  \Updownarrow \\
  \operatorname{Supp} E^{r-i} \subset Z \\
  \Updownarrow \\
  H^i_x(F_x) = 0 \text{ pour tt } x \text{ fermé dans } U.
\end{array}
\right.
\]
\note{« dans $U$ » est souligné.}

\[
\left.
\begin{array}{l}
  \alpha_i \text{ bijectif pour } m \text{ grand},\ i \leq n \\
  \Updownarrow \\
  \operatorname{Supp} E^{r-i} \subset Z \text{ pour } i \leq n \\
  \Updownarrow \\
  \operatorname{prof}(F_x) > n \text{ pour tt } x \text{ fermé dans } U \\
  \Updownarrow \\
  H^i(U, F(-m)) \leftarrow H^i(X, F(-m)) \text{ pour } m \text{ grand},\ i < n \text{ et} \\
  H^n(U, F(-m)) \to H^{n+1}_Z(X, F(-m)) \text{ injectif}
\end{array}
\right.
\]
\note{L'avant-dernière ligne est surchargée : le premier membre est récrit ($U$ sur $X$) ; le second, « $H^i(X, F(-m))$ injectif pour $i \leq n$ », est barré d'un trait, et « pour $m$ grand, $i < n$ et » est récrit au-dessus ; la flèche est un double trait. À la dernière ligne, « surjectif pour $i \leq n$ » est biffé et remplacé par « injectif ».}

\[
(\mathrm{B}_n) \quad
\left.
\begin{array}{l}
  \alpha_i \text{ bijectif pour } m \text{ grand},\ i \leq n, \\
  \quad \text{et } \alpha_i \text{ injectif pour } m \text{ grand},\ i = n+1 \\
  \Updownarrow \\
  \operatorname{Supp} E^{r-i} \subset Z \text{ pour } i \leq n,\ \operatorname{Supp} E^{r-n-1} \ldots Z \\
  \Downarrow \\
  \operatorname{prof}(F_x) > n \text{ si } x \text{ fermé dans } U, \text{ et} \\
  \operatorname{prof}(F_x) > n+1 \text{ si de plus } x \in V \cap U,\ V \text{ voisinage ouvert convenable de } Z \\
  \Updownarrow \\
  H^i(U, F(-m)) = 0 \text{ pour } m \text{ grand},\ i \leq n \\
  \Updownarrow \\
  \coprod_{m \geq 0} H^i(U, F(m)) \text{ de type fini sur } S \text{ pour } i \leq n
\end{array}
\right.
\]
\marginal{$H^i(U, F)$ de dim finie pour $i \leq n$ (cf 1°)}
\note{Le label « $\mathrm{B}_n$ » est encerclé en marge gauche. La fin de la quatrième ligne, « $\operatorname{Supp} E^{r-n-1}$ \ill{} \ill{} \ill{} $Z$ pour \ill{} \ill{} \ill{} », ne se lit pas (on la rend par des points) ; un « prof » est biffé en tête de ligne. Après « de type fini sur $S$ », un signe surchargé (« $(-h)$ » ?) ; lecture incertaine. La note marginale, encerclée, est reliée à l'accolade par une double flèche marquée « cf 1° ».}

Cela résout la question (\uncertain{a}). Si, sous les conditions de $(*)$, on a donc [si $t$ est $F$-régulier]
\[
\left\lbrace
\begin{array}{l}
  H^i(U, F) \simeq H^i(\hat{X}, \hat{F}) \quad \text{est} \quad
  \left\lbrace
  \begin{array}{l}
    \text{isom si } i < n \\
    \text{mono si } i = n
  \end{array}
  \right. \\[1ex]
  H^i(\hat{X}, \hat{F}) \simeq \varprojlim H^i(X_m, F_m) \text{ pour } i \leq n
\end{array}
\right.
\]

\page{58}
ce qui donne des conditions suffisantes pour la validité de c). Cependant, elles ne sont pas nécessaires [\ill{} on exigerait la validité de c pour $i < n$, et tous les tordus de $F$], comme on le voit en faisant $n = 1$ : alors il suffit que on ait $\operatorname{prof} F_x \geq 2$ pour $x$ fermé dans $U$, pour pouvoir conclure que
\[
  H^0(U, F) \longrightarrow H^0(\hat{X}, \hat{F})
\]
est bijectif.

(2°) La suite exacte ci-dessus montre que l'homom.
\[
  H^i(U, F) \longrightarrow H^{i+1}_Z(X, F)
\]
est un isom.\ \uncertain{modulo} \uncertain{vectoriels} finis. \struck{Donc} \struck{$H^i(U, F)$}
\[
(\mathrm{A}'_i) \quad
\left.
\begin{array}{l}
  H^i(U, F) \text{ de dim finie sur } k \\
  \Updownarrow \\
  H^{i+1}_z(F_z) \text{ de dim finie [i.e.\ de long.\ finie] pour tt } z \in Z \\
  \Updownarrow \\
  z \in Z \Rightarrow z \text{ pt isolé de } \operatorname{Supp} E^{r-i-1} \cup \lbrace z \rbrace \\
  \Updownarrow \\
  \alpha_{i+1} \text{ injectif pour } m \text{ grand} \\
  \Updownarrow \\
  H^i(X, F(-m)) \to H^i(U, F(-m)) \text{ surjectif pour } m \text{ grand}
\end{array}
\right.
\]
\note{Le label « $\mathrm{A}'_i$ » est encerclé en marge gauche. À la troisième ligne, « $\in$ » est récrit sur « $\notin$ » ; à la dernière, le $U$ est récrit sur un $X$.}

\page{59}
$R^p f_* [R^q j_* (F(m))]$
\note{En haut à droite ; l'argument de $R^q j_*$ est surchargé et d'une lecture incertaine.}

\struck{$H^i(U, F(-m)) \to H^i_Y$ \ldots}

\struck{$H^i(U, F(m))$}
\note{Le haut de la page, au-dessus d'un trait horizontal, est annulé par un grand trait oblique ; de légers traits obliques au crayon traversent aussi le paragraphe suivant.}

(3°)
\[
\left.
\begin{array}{l}
  \coprod_{m \geq 0} H^i(U, F(m)) \text{ de type fini sur } S \text{ pour } i \leq n \\
  \Updownarrow \\
  R^i g_*(F) \text{ coh.\ pour } i \leq n, \text{ i.e.\ } \underline{H}^i_Z(F) \text{ coh.\ pour } i \leq n+1, \\
  \quad \text{i.e.\ } H^i_z(F_z) \text{ de dim finie pour } i \leq n+1,\ z \in Z
\end{array}
\right.
\]

\[
(\mathrm{A}^f_n) \quad
\left.
\begin{array}{l}
  H^i(U, F) \text{ de dim finie si } i \leq n \\
  \Updownarrow \\
  F \text{ de prof} > n+1 \text{ sur } U \text{ au voisinage de } Z \\
  \Updownarrow \\
  z \in Z \Rightarrow z \text{ pt isolé de } \operatorname{Supp} E^{r-i-1} \cup \lbrace z \rbrace \text{ si } i \leq n \\
  \Updownarrow \\
  \alpha_i \text{ injectif pour } \underline{m \text{ grand}} \text{ si } i \leq n+1 \\
  \Updownarrow \\
  H^i(X, F(-m)) \to H^i(U, F(-m)) \text{ surjectif pour } m \text{ grand},\ i \leq n \\
  \Updownarrow \\
  R^i g_*(F) \text{ coh.\ pour } i \leq n \iff \underline{H}^i_z(F) \text{ de dim finie pour } z \in Z,\ i \leq n+1 \\
  \Updownarrow \\
  \coprod_{m \geq 0} H^i(U, F(m)) \text{ de type fini sur } S \text{ pour } i \leq n
\end{array}
\right.
\]
\note{Le label « $\mathrm{A}^f_n$ » est encerclé en marge gauche ; l'exposant, lu $f$, est incertain. Une double flèche de la page n'a qu'une pointe (vers le haut) ; on l'a rendue $\Updownarrow$ comme les autres.}

\page{60}
\note{Page dactylographiée, tournée tête-bêche : un second exemplaire du feuillet de la page 56 (même fin d'énoncé, « Corollaire~3.2. Sous les conditions de 3.1. ( »). Ce qui suit est de sa main, au crayon.}

\[
  \sum_{p \geq 0} R^i f_{0*} F_0(q - p) \qquad\qquad q - p \geq n
\]
\[
  S_{pq} = \left\lbrace
  \begin{array}{ll}
    0 & \text{si } p > 0 \\
    S_q & \text{si } p = 0
  \end{array}
  \right.
\]
\note{L'indice de $S$ est surchargé (lu $pq$, lecture incertaine). Suivent deux croquis au crayon : un quadrant d'axes $p$ (horizontal) et $q$ (vertical), avec une bande étroite le long de l'axe des $p$ et la légende « $F_0(-p)$ » ; à droite, un second quadrant traversé d'une diagonale.}

\end{document}
